\begin{abstract}
Power outages can disable O-RAN radio units (RUs) as their finite local batteries deplete, threatening emergency connectivity when communication is most needed. Operators therefore need to conserve energy without causing service to fall below an emergency target. This study asks how predefined RU sleep policies affect emergency connectivity and RU battery lifetime during a network-wide outage and whether they meet a served-user target of 0.70. It assesses an Always Active (AA) baseline against Staggered Active (SA), which alternates fixed RU groups, and Threshold Staggered Active (TSA), which postpones staggering until all RUs reach a battery threshold. The evaluation compares served-user fraction, horizon-average emergency QoS, mean RU-depletion time, and SLA-based network lifetime within a simplified O-RAN outage model. AA provides the strongest service baseline, with a median SLA-based network lifetime of 21.0 steps. SA lengthens median mean RU-depletion time from 31.2 to 46.2 steps (48\%) but has a median network lifetime of zero. TSA retains near-AA network lifetime (20.5 steps) while extending median mean RU-depletion time to 33.4 steps. Thus, among the fixed policies TSA best preserves target service but modestly extends battery availability, indicating adaptive outage management is needed to make material energy savings without jeopardizing emergency connectivity.
\end{abstract}

\begin{IEEEkeywords}
O-RAN, radio unit, sleep policy, power outage, simulation
\end{IEEEkeywords}
